\section{Introduction}
Team contests test more than whether a model can solve an isolated problem.
Contestants must divide a packet, allocate scarce time and tools, exchange
partial results, verify work, and decide when to submit. Agent Olympiad
therefore evaluates \emph{teams under contest constraints}, not unconstrained
solver ability. A strong solo system may remain a useful baseline, but it does
not answer whether a group forms a coherent team.

We use four ability families. \textbf{Problem solving and decomposition}
cover solving and splitting the task. \textbf{Strategy and resource
allocation} cover triage, assignment, timing, tool use, and submission policy.
\textbf{Communication} covers useful transfer of questions, evidence, and
observations. \textbf{Cohesion} asks whether interactive teammates improve on
an isolated orchestrator with independent workers, rather than merely
parallelizing calls.

The intended contributions are: (1) competition-level tasks paired with
source-grounded rule profiles; (2) a baseline matrix that separates extra
sampling, decomposition, communication, architecture, and budget; and (3) an
evaluation stack that reports task scores beside collaboration measures and
explicitly marks proxies and missing comparisons. Current evidence is
preliminary. We do not claim that teams generally outperform solo agents, nor
that turn limits alone establish harm from collaboration.
